\section{A bounded kernel in the homogeneous multiplicity branch}
\label{sec:heisenberg-homogeneous-multiplicity}

\paragraph{Status and provenance.}
D-TD27--29, reserved question Q2en, is a derivation by
\texttt{/root/imprimitive} following D-TD20--21.  Root proposed
examining the modular multiplicity action after homogeneity; the
finite-field factorization, quotient-power estimate and bounded-kernel
central-product reduction below are supplied here.  Complete author
proof, source-interface, checker and first-output readings pass;
root and independent whole-proof reviews pass.  The irreducible-kernel case and the symmetric
sign cases are deliberately outside the new conclusion.

Let \(\mathcal Q\) be a nonzero symplectic space over \(\mathbf F_p\),
and let \(H\le\operatorname{Sp}(\mathcal Q)\) act irreducibly over
that original field.  Suppose that there is an actual surjection
\[
 \psi:H\longrightarrow A_n,\qquad n\ge6,\qquad
 N=\ker\psi,\qquad
 \mathcal Q\big|_N\simeq S^{\oplus r},
 \qquad d=\dim_{\mathbf F_p}S.
 \tag{TDM.1}
\]
Here \(S\) is irreducible over \(\mathbf F_p\).  This is the
alternating, homogeneous branch of D-TD21.  It is not a statement
about the characteristic-zero kernel of an entire standard partner.
The group H is its actual faithful matrix image; in particular its
subgroups have their actual actions on \(\mathcal Q\).

\subsection*{D-TD27. The finite-field multiplicity map and geometric rank}

\begin{lemma}[D-TD27]
There are integers \(f,a\ge1\), with \(d=af\), and an intrinsic
field \(F=\mathbf F_{p^f}\), for which
\[
 \mathbf F_p[N]\big|_{\mathcal Q}
       \simeq\operatorname{Mat}_a(F)\otimes1_r,
 \qquad
 \mathcal Q\simeq F^a\otimes_F F^r
       \quad\hbox{as an underlying }\mathbf F_p\hbox{-space}.
 \tag{TDM.2}
\]
H fixes F pointwise and every element has a factorization
\(h=A_h\otimes B_h\) over F, unique up to reciprocal nonzero
scalars.  The maps \([A_h]\) and \([B_h]\) are actual projective
homomorphisms.  The second one kills N.  If \(r>1\), it induces
a faithful projective representation
\[
           A_n\lhook\joinrel\longrightarrow\operatorname{PGL}_r(F),
 \qquad
 h\notin N\ \Longrightarrow\
          \operatorname{rank}_{\mathbf F_p}(h-1)\ge d.
 \tag{TDM.3}
\]
No semisimplicity of an individual h, no absolute irreducibility of
S over \(\mathbf F_p\), and no splitting of \(H\to A_n\) is assumed.
\end{lemma}

\begin{proof}
Let \(A\) be the image of \(\mathbf F_p[N]\) on S.  It acts
faithfully and irreducibly.  Its nilpotent radical annihilates S:
the radical times S is either zero or all of S, and the latter
alternative contradicts nilpotence.  Faithfulness makes the radical
zero.  Artin--Wedderburn therefore identifies A with one full matrix
algebra over a division ring, since a faithful simple module cannot
be supported on only one of several nonzero simple factors.
The division ring is finite, hence a field F.  The unique simple
module and its commuting algebra identify S with \(F^a\) and
\(\operatorname{End}_N(S)=F\), with the harmless opposite-ring
convention removed by commutativity.  These are the classical
finite-dimensional statements in
\cite[ChapterII, Theorem2.2, pp.~83--84; Theorem2.13, pp.~91--92;
Theorem2.48, pp.~117--118]{TD:Knapp2016}.
The homogeneous restriction gives the actual algebra and tensor
identification in (TDM.2), with F the center of this algebra.

Normality of N makes H normalize A and its center F.  Conjugation
by N acts trivially on that center.  The induced field-automorphism
action thus factors through \(A_n\).  The automorphisms of F over
\(\mathbf F_p\) form the cyclic group generated by Frobenius:
its f distinct powers already exhaust the at most f embeddings of
this degree-f field extension.  As \(A_n\) is perfect, its action
on that cyclic group is trivial.  H is therefore F-linear.

An F-linear automorphism of \(\operatorname{Mat}_a(F)\) is inner,
by Skolem--Noether
\cite[Theorem2.41 and Corollary2.42, pp.~113--114]{TD:Knapp2016}.
Choose an implementer \(A_h\).  Dividing h by
\(A_h\otimes1\) gives an element of the commuting algebra
\(1\otimes\operatorname{Mat}_r(F)\), hence \(1\otimes B_h\).
The two implementers can change only by reciprocal scalars;
composition consequently makes their projective classes genuine
homomorphisms.  Elements of N have trivial second projective class.

If the resulting map \(A_n\to\operatorname{PGL}_r(F)\) were
trivial and \(r>1\), every subspace \(F^a\otimes L\), for an
F-line \(L\subset F^r\), would be H-invariant.  This is a
nonzero proper \(\mathbf F_p\)-subspace, contrary to irreducibility.
The map is nontrivial, and simplicity of \(A_n\) makes it faithful.

We give the geometric tensor-rank argument explicitly because
finite-order operators in characteristic p need not be semisimple.
For arbitrary invertible matrices \(A_0,B_0\) of respective sizes
a,r over a field, let \(\nu_{\rm alg}(B_0)\) be r minus its
largest geometric eigenspace dimension over an algebraic closure.
For every scalar \(\lambda\) in that closure,
\[
 \dim\ker(A_0\otimes B_0-\lambda)
       \le a\bigl(r-\nu_{\rm alg}(B_0)\bigr),\qquad
 \operatorname{rank}(A_0\otimes B_0-\lambda)
       \ge a\nu_{\rm alg}(B_0).
 \tag{TDM.4}
\]
Indeed, identify the kernel with intertwining maps between the
operator \(\lambda A_0^{-\mathsf T}\) and \(B_0\).
A Jordan block of size u in the first operator and a block of
size v in the second contribute zero for different eigenvalues
and \(\min(u,v)\) otherwise.  For a fixed first block their total
contribution is at most u times the number of second blocks at
the matching eigenvalue.  That number is the corresponding
geometric eigenspace dimension.  Summing u over the first blocks
proves (TDM.4) in every characteristic.

For \(h\notin N\), faithfulness of its second projective image
gives \(\nu_{\rm alg}(B_h)\ge1\).  Apply (TDM.4) at
\(\lambda=1\), then multiply the F-rank by f.  This gives
\(\operatorname{rank}_{\mathbf F_p}(h-1)\ge fa=d\), as claimed.
\end{proof}

\subsection*{D-TD28. The actual quotient-power budget}

Now put \(\dim_{\mathbf F_p}\mathcal Q=2m\).  Suppose that
\(P\) is an actual finite projective factor image on a complex
space W of dimension \(p^m\), with a nondegenerate projective
Heisenberg normal subgroup identified with \(\mathcal Q\), and
\(P/\mathcal Q=H\) with the conjugation action above.
Suppose an actual sphere source has a factorwise action
\[
 J\twoheadrightarrow P\twoheadrightarrow H
       \mathop{\twoheadrightarrow}^{\psi} A_n,
 \qquad V=W\otimes R,\quad e=\dim V,
 \qquad
 \sum_i\operatorname{rank}(1-V(g_i))=2e+6.
 \tag{TDM.5}
\]
The tuple \((g_i)\) generates J and has product one.  R is the
entire external partner; no count of genuine projected marks is
imposed.  The individual factor actions may be projective, while
the full coefficient action is the actual finite linear action.

\begin{theorem}[D-TD28]
Under (TDM.1)--(TDM.5), if \(r>1\), then
\[
             e\bigl(1-85p^{-d/2}\bigr)\le252.
 \tag{TDM.6}
\]
In particular, if \(e>504\), then
\[
 p^d<28900,
 \qquad d\le14,
 \qquad |N|<\mathcal B:=28900^{14}.
 \tag{TDM.7}
\]
These bounds concern the kernel module and the kernel group, not
the homogeneous multiplicity r or the Heisenberg half-rank m.
\end{theorem}

\begin{proof}
The exact Weyl-basis trace identity of D-TD05 gives, for any
projective normalizer lift T of \(h\in H\), in unitary normalization,
\[
 \frac{|\operatorname{tr}T|}{p^m}
   \le p^{-\operatorname{rank}(h-1)/2}.
 \tag{TDM.8}
\]
The trace may be zero; that only strengthens the bound.  This
identity includes all Heisenberg translations and all scalar
phases, also at p=2.  If \(\psi(h)\ne1\), D-TD27 therefore
bounds (TDM.8) by \(c=p^{-d/2}\).
Tensoring the entire partner cannot increase normalized absolute
trace.  Compatible unitary normalizations are available for the
finite projective factor images, and the remaining scalar in the
finite full operator has absolute value one.  Consequently
\( |\operatorname{tr}V(g)|/e\le c\) for every actual g whose
\(A_n\)-image is nonidentity.

The power calculation must use an actual lift period, since W need
not be a projective representation of \(A_n\) itself.  Write L
for the order of g in J, and o for the order of its image in
\(A_n\).  Thus \(o\mid L\) and \(V(g)^L=1\); Fourier inversion
with this period is valid even if the operator has smaller order.
Among \(0\le j<L\), precisely \(L/o\)
powers have trivial quotient image.  Bound their normalized
absolute traces by one and every other power by c.  Fourier
inversion, for each target eigenvalue of the full operator, gives
\[
 \frac{\operatorname{rank}(1-V(g))}{e}
   \ge (1-c)(1-1/o).
 \tag{TDM.9}
\]
This reasoning retains arbitrary lift orders and kernel powers;
it never replaces L by o for the full operator.

The actual \(A_n\)-quotient of the given connected sphere cover
is connected.  Since \(n\ge6\), its group is nonsolvable and
has order greater than sixty.  The complete genus-zero and
genus-one exclusions and the elementary hyperbolic signature
calculation in D-TD02 imply
\(\sum_i(1-1/o_i)\ge85/42\), with identity images omitted.
The same conclusion follows directly from that accepted proof,
without treating W as an \(A_n\)-module.  Summing (TDM.9) yields
\(2+6/e\ge(1-c)85/42\), which is exactly (TDM.6).

If \(p^d\ge170^2=28900\), then \(c\le1/170\), so
(TDM.6) forces \(e\le504\).  This contradicts the stated
strict opposite assumption.  Because \(2^{15}>28900\), the
first inequality in (TDM.7) gives \(d\le14\).
N acts faithfully on S: an element acting trivially there acts
trivially on all of \(S^{\oplus r}=\mathcal Q\), hence is
the identity in the actual subgroup H.  Therefore
\[
 |N|\le|\operatorname{GL}_d(p)|<p^{d^2}
                =(p^d)^d<28900^d\le\mathcal B.
 \tag{TDM.10}
\]
This proves every asserted bound without restricting r.
\end{proof}

\subsection*{D-TD29. An actual quasisimple lift and its defining field}

\begin{theorem}[D-TD29]
Retain D-TD28 with \(r>1\), \(e>504\), and
\(n\ge\mathcal B+2\).  There is an actual normal quasisimple
subgroup \(E\lhd H\) for which
\[
 H=NE,\qquad [N,E]=1,\qquad
 N\cap E=Z(E),\qquad E/Z(E)\simeq A_n.
 \tag{TDM.11}
\]
This is an actual central product in the defining symplectic
matrix group, not a lift assumed from an abstract quotient.

More precisely put \(F_0=\operatorname{End}_H(\mathcal Q)\).
It is a field \(\mathbf F_{p^{f_0}}\), containing F, and one
can write
\[
 \mathcal Q\simeq F_0^a\otimes_{F_0}F_0^s,
 \qquad r=[F_0:F]s,
 \qquad
 N\subset\operatorname{GL}_a(F)\otimes1,
 \quad E\subset1\otimes\operatorname{GL}_s(F_0).
 \tag{TDM.12}
\]
The E-action on \(F_0^s\) is faithful and absolutely irreducible,
and its projective kernel is exactly \(Z(E)\).  In this realization
\[
 p^{a f_0}<28900,
 \qquad a f_0\le14,
 \qquad |F_0|<28900.
 \tag{TDM.13}
\]
No upper bound on s, r, m or n is concluded.
\end{theorem}

\begin{proof}
Conjugation defines a homomorphism \(H/N\to\operatorname{Out}(N)\).
Every automorphism of the finite group N permutes its underlying
set, so \(|\operatorname{Out}(N)|\le|N|!\le\mathcal B!\).
If the displayed homomorphism from the simple \(A_n\) were
nontrivial, it would be injective.  This is impossible because
\(n!/2>\mathcal B!\) for \(n\ge\mathcal B+2\).
Thus every conjugation action of H on N is inner.

Put \(C=C_H(N)\).  Multiplying an arbitrary element of H by
an appropriate element of N makes it centralize N, so \(H=NC\).
Moreover \(N\cap C=Z(N)\), and \(C/Z(N)\simeq A_n\).
Set \(E=C'\).  Perfectness of \(A_n\) implies
\(C=E Z(N)\).  Since \(Z(N)\) centralizes C,
\(E=[C,C]=[E,E]\); thus E is perfect and maps onto \(A_n\).
Its central kernel is \(E\cap Z(N)\), and an element of
\(Z(E)\) must map to the trivial center of \(A_n\).
Consequently \(Z(E)=E\cap Z(N)=E\cap N\), proving (TDM.11).
C is normal in H, and E is characteristic in C, so E is normal
in H as asserted.

The exact commuting algebra in (TDM.2) identifies
\(C\subset1\otimes\operatorname{GL}_r(F)\).
Thus E acts faithfully on the multiplicity space \(M=F^r\).
That action is irreducible over F: an invariant proper F-subspace
of M would, after tensoring with \(F^a\), give an H-invariant
proper \(\mathbf F_p\)-subspace of \(\mathcal Q\), since H=NE.
Apply the same finite-dimensional algebra theorem as in D-TD27
to the image of F[E] on M.  It is
\(\operatorname{Mat}_s(F_0)\) for a finite extension \(F_0/F\),
and M is its unique simple module.  In particular
\(\operatorname{End}_E(M)=F_0\).
Commuting first with N and then with E shows that this field is
exactly \(\operatorname{End}_H(\mathcal Q)\).
This proves (TDM.12), including \(r=[F_0:F]s\).

Over \(F_0\) the E-matrices span the full matrix algebra, which
remains a full matrix algebra after any field extension.  Their
representation is therefore absolutely irreducible.  It is still
faithful.  A scalar E-matrix commutes with every E-matrix and hence
comes from \(Z(E)\); conversely central elements act as scalars
in this absolute model.  The projective kernel is precisely that
center.

Write \(h=uz\), with \(u\in N,z\in E\), using (TDM.11).
If \(\psi(h)\ne1\), then \(z\notin Z(E)\); its multiplicity matrix over \(F_0\) is
nonscalar.  Formula (TDM.4) over \(F_0\) now gives
\(\operatorname{rank}_{\mathbf F_p}(h-1)\ge a f_0\).
Repeat exactly the full-order Fourier and actual-cover calculation
of D-TD28 with d replaced by \(a f_0\).  Since \(e=\dim V>504\),
it follows that \(p^{a f_0}<28900\).  The remaining inequalities
in (TDM.13) follow immediately.
\end{proof}

\paragraph{Application threshold and remaining boundary.}
In the maximal-standard application of D-TD19, the previously fixed
threshold \(N_*\) exceeds \(\mathcal B+2\).  For example its
definition has \(N_*>240!\), while
\(240!\ge120^{120}>28900^{30}>28900^{14}+2\), using
\(120^4>28900\).  The full source degree is then certainly
greater than 504.  Thus D-TD29 applies throughout that application
whenever its homogeneous multiplicity r exceeds one.

The case \(r=1\) is not disposed of by a trivial multiplicity
map: N can then be irreducible on the entire defining space, and
the quotient can act through its normalizer.  The \(S_n\) quotient
also has potential sign field automorphisms or two exchanged
isotypic components, and is not included in D-TD27--29.  Finally,
the bounded finite field in (TDM.13) is a defining field for a
modular E-representation.  It is unrelated to equality of the
original characteristic-zero coefficient field or to a choice of
polarization.  No source existence, complete symplectic-normalizer
replacement, or rank comparison with the natural n-letter module
has been inferred.
